\section{核心代码}

本文核心算法使用Python 3.10实现，依赖库包括numpy 1.24、pandas 2.0、scikit-learn 1.3、shap 0.42等。完整代码已提交至附件文件夹，此处展示关键模块。

\subsection{问题一：多模态特征对齐}

\begin{Python}{统计特征提取与对齐}
import numpy as np
import pandas as pd

def extract_statistical_features(X_seq):
    """从变长时序特征提取统计量"""
    features = []
    for seq in X_seq:
        if len(seq) > 0:
            features.append([
                np.mean(seq, axis=0),
                np.std(seq, axis=0),
                np.max(seq, axis=0),
                np.min(seq, axis=0),
                np.percentile(seq, 25, axis=0),
                np.percentile(seq, 75, axis=0)
            ])
        else:
            features.append(np.zeros((6, seq.shape[1])))
    return np.concatenate([f.flatten() for f in features])
\end{Python}

\subsection{问题二：鲁棒梯度提升模型}

\begin{Python}{模态缺失训练策略}
from sklearn.ensemble import GradientBoostingRegressor
import numpy as np

class RobustGBDT:
    def __init__(self, dropout_rate=0.15):
        self.model = GradientBoostingRegressor(
            n_estimators=200,
            learning_rate=0.05,
            max_depth=7,
            random_state=42
        )
        self.dropout_rate = dropout_rate

    def apply_modal_dropout(self, X, modal_dims):
        """训练时随机丢弃模态"""
        X_aug = X.copy()
        for i in range(len(X)):
            for start, end in modal_dims:
                if np.random.rand() < self.dropout_rate:
                    X_aug[i, start:end] = 0
        return X_aug
\end{Python}

\subsection{问题三：SHAP可解释性分析}

\begin{Python}{SHAP值计算与可视化}
import shap

# 训练完成后计算SHAP值
explainer = shap.TreeExplainer(model)
shap_values = explainer.shap_values(X_test)

# 全局特征重要性
shap.summary_plot(shap_values, X_test,
                  feature_names=feature_names,
                  show=False)

# 单样本解释
shap.force_plot(explainer.expected_value,
                shap_values[idx],
                X_test[idx],
                feature_names=feature_names)
\end{Python}
